\documentclass{article}
\usepackage{amsmath}
\usepackage{tikz}
\usepackage{tkz-euclide}
\usepackage{amsfonts}
\usepackage{graphicx}
\graphicspath{{./assets/}}

\numberwithin{equation}{section}
\title{Generell matematikk notat}
\author{Martin N. Håvardsen}
\date{\today}

\begin{document}
\maketitle

\section{Abel oppgave 1}
Finn det største heltallet $k$ der det finnes et heltall $a$ som oppfyller:
\begin{align}
  N=\displaystyle\frac{(a+k)!}{a !} \textbf{ er \textit{IKKE} delelig på 2026 }
\end{align}
Vi vet at $k<2026$. Dersom $k=2026$ så må $N$ har faktorer som fyller ut intervallet $\langle a, a+k]$. En faktor delelig på $2026$ vil oppstå \textbf{minst} en gang hver $2026$ heltall. Dvs. $2026,\quad 2026\cdot 2,\quad 2026\cdot 3, \quad (\dots)$\;. Intervallet må derfor inneholde et heltall delelig på $2026$. \newline
Er det flere slike tilfeller? Ingen heltall opptil 2026 vil være delelig på 2026. Og det har vist seg at $k<2026$, så svaret må derfor være $k=2025$. Men dette er for lett. Det er en sjanse for at jeg har skrevet opp oppgaven feil (jeg har ikke den fremme her nå). Uansett så må svaret for oppgaven oppgitt her være: $k=2025$
\section{Calculus fremgang}
Calculus har fremstått som noe vanskelig og krevende. Tvert imot, det er ikke noe særlig vanskelig. \textit{Calculus boken} går bare gjennom noen få flere ting enn R-matten gjør. Det er liten forskjell. Den eneste forskjellen er hvor fokuset ligger, og hvordan konsept blir introdusert. \newline
\textit{Calculus boken} går gjennom legger frem tankeganger for en hypotetisk oppdaging av konseptene. Dvs. den legger frem nye idéer, deretter brukes gammel teori for å utlede egenskaper etc. knyttet til idéene. For eksempel, [...] \newline
R-matten introduserer de nye temaene på kort vis, følgende kommer formler. Det blir lagt vekt på applikasjonene til disse formlene. Sjeldent får man bevis eller forklaring; sterk kontrast til \textit{Calculus boken}. \newline
Jeg har en mistanke om at \textit{Linear Algebra boken} er mer omfattende og krevende enn \textit{Calculus boken}. Det er vanskeligere å komme seg gjennom den (i alle fall i min erfaring). Kanskje konseptene er mer unike og motstridende gamle oppfatninger. Oppgavene i \textit{Linear Algebra boken}, sånn jeg husker dem, har to arter: enten gratis (menneskekalkulator-oppgaver) eller komplekse (blanding av teori med unik løsningsteknikker). Dette kan være farget av mitt da unge sinn (xd).\newline
\section{Skriving og personlig ytring}
Jeg liker ikke å skrive. Jeg tror kanskje at det er på grunn av at jeg må skrive på norsk (kun). Vanligvis så snakker jeg en norsk-engelsk hybrid rikt med moderne internettkultur slang. Jeg har opplevd tilfeller der jeg kun fant passende \textbf{engelske} ord, selv om jeg talet på norsk. Det er mulig at dette gjør det mer utfordrende å uttrykke mer presise idéer. \newline
Det føles også avgrensende ut. Det er vanskelige å formulere tankene mine når jeg må overvåke egen rettskriving, ordbruk, struktur. Jeg kan bare skrive så så mye. Resultatet må derfor være mer konsentrert og konsist; det er vanskelig å skrive. Spesielt dersom verket ikke skal dreie om noe særlig verdifullt. Da blir det bare en hel mye bortkastet skrivetid. \newline
Det er kanskje mulig å fikse dette med å senke kvalitetsstandarden, tillate engelsk ordbruk, tillate lengre tekstbidrag. Men det er ikke så lett. Kodeveksling er ikke passende for mer formelle og faglige kontekster. Jeg vil ikke skrive så mye uformellt \textit{piss} xd. Lengre tekster blir ikke lest. Jeg gidder ikke å lese en lang, uformell, selv-skreven tekst. En lang tekst vil derfor ikke bli lest, så den vil i praksis ikke fungere som lagring av data eller informasjon. Bare tenkeskrivingstekster bør være lange; tekster med lite innhold. Ironisk. Å senke kvalitetsstandarden er lovende, men ubehagelig. Jeg klarer som regel å oppdage feil/svakheter i skrivingen min. Det å rette dem er tidkrevende og distraherende. Men det er i seg selv distraherende å la dem være.\newline
Jeg har ingen gode svar her. Jeg skal prøve mitt beste på å overse feil/svakheter i tekstene mine. Kanskje det vil pådrive en økning i litteratursproduksjonen min \textit{xd idk mbe ig}. Det er kanskje fornuftig å spare skriving på de mer innholdsrike, komplekse temaene. Derimot så er tenkeskriving en profesjonelt validert teknikk. De to skrivestrategiene må separeres, siden formålene er såpass ulike. \newline
\tiny{Ble inspirert av \textit{Flowers for Algernon}. :)}
\small
\section{Hva er den beste måten å forstørre norsk vokabulær?}
Å forstørre vokabulæret innebærer hovudsakelig bare å huske flere ord. Med et stort mangfold av ord kan man uttrykke seg mer presist og fargerikt. Dette hjelper særlig i akademisk kontekst. For eksempel i matematikk, norsk, fysikk, historie (fag som jeg har nå). \newline
Det finnes allerede mange teknikker for å pugge nytt innhold. \textit{Flashcards} er et fundament i denne verdenen. De åpner for enkel jevn øving med varig effekt. Problemt ligger ikke i kvaliteten av resultatene, men heller kostnadene av prosessen. Å lage selve kortene er tidkrevende. Å finne ord som det er "verdt" å legge til i kortstokken er i seg selv en oppgave. Kvaniteten denne teknikker tilbyr er altså for dårlig og passer bedre for mer konsentrert informasjon; ikke tilfeldige, diskontinuerlige datapunkt med kort, overflatisk innhold. \newline
Videre har vi: konsumpsjon av norsk litteratur, spesielt sakprosa. Her gjør underbevisstheten alt arbeidet. Teknikken dreier seg ikke om å utsøke nye ord, men heller arbitrært støte på dem. Kanskje via lesing, eller kanskje lytting av lydbøker. Målet er fremdeles å fange inn nye ord, men selve prosessen vil forbli en smule "morsomt". Konvensjonell bruk av ny vokabulær vil også være bygd inn i læreprosessen. Ikke bare får du lære ordenes betydning, men også klassisk bruk og personlig assosiasjon med ordet. \newline 
Begge metodene er ineffektive. Rent talt så gidder jeg ikke å prokurere flere ord til kortstokkene mine. Det er heller ikke noe særlig gøy å bruke kortstokkene xd. På den andre siden så er jeg ikke så interessert i norsk litteratur. Jeg ønsker heller å lese engelsk-til-norsk oversettelser av mer berømte verk (større sjanse for innlevende fortellinger :D). Dessuten så gidder jeg ikke å sjekke opp eventuelle nye ord som dukker opp. Det dra meg ut av lesing tross alt. Dette har ført til at jeg ikke har gjort noe som helst for å forbedre ordforrådet. For meg så er det enten eller. Resultatene må være fenomenale eller så må prosessen vær for den simple å utføre. Ingen av disse er tilfellet :(
\end{document}

